\documentclass[a4paper,12pt]{article}
\usepackage{amsmath,amssymb}
\usepackage{graphicx}
\usepackage[hidelinks]{hyperref}
\usepackage[margin=2.5cm]{geometry}
\usepackage[numbers]{natbib}
\usepackage{booktabs,array,float}
\title{From Compute to Complexity: A Physical Theory of Intelligence Emergence}
\author{Qinrang Liu}
\date{June 2026}
\begin{document}
\maketitle

Theory of Intelligence Emergence and Its Implications for Artificial General Intelligence Abstract The rapid scaling of large language models has delivered remarkable functional capabilities yet produced exponentially growing energy costs with sub-linear returns---a thermodynamic trajectory that converges not toward general intelligence but toward an unsustainable asymptote. We argue that this trajectory is not an engineering deficiency but a consequence of pursuing the wrong variable: compute,
rather than complexity. Von Neumann identified in 1948 that intelligence requires a complexity threshold; here we quantify that threshold through a framework grounded in thermodynamic phase transitions, renormalization group theory, and complex network science. The result is the Coordination Spatiotemporal Complexity theorem: CST = ( Sc$\cdot$Tc)$\cdot$exp($\alpha$$\cdot$$\Gamma$st), where structural integration, dynamical richness, and their physical coupling jointly de- termine emergent
intelligence potential. We derive six universal thresholds at natural constants {1/$\sqrt{}$ 2,1, $\phi$, e, $\pi$, $\delta$ }and validate across 40 biological and artificial systems spanning 8 taxonomic grades and 18 distinct ANN/NMH architectural families (Spear- man$\rho$= 0.976, 100% accuracy under UCCP normalization). Neuromorphic hardware (Intel Loihi-2) is separately classified from binary-digital ANN, confirming the $\alpha$-barrier

prediction. Intelligence Efficiency $\eta$Ireveals an approximately six-order-of-magnitude gap between brains and current AI, and a four-generation hardware roadmap identifies the physically necessary path from present systems to general intelligence. Keywords: intelligence emergence; complexity threshold; von Neumann; spatiotem- poral coordination; intelligence efficiency; phase transitions; neuromorphic computing 2

\end{document}